%% ==========================================================================
\section{Empirical grounding in the registry}
\label{sec:registry}
%% ==========================================================================


The two stipulations that lean on the physical-attack registry~\cite{lopp-attacks}
--- the coercion-time budget $t$ (the limitations section in~\cite{deadly-race}) and the murder-cost weighing of the
residual-window analysis (the residual deadly-race window in~\cite{deadly-race}) --- admit directional empirical support. The registry records each
incident only as \emph{date, victim, location, headline}; it carries no duration or
outcome field, so we code the $352$ incident headlines (at commit \texttt{9a4a62a}) by
keyword. The coding is coarse and the categories overlap; a reproducible script
accompanies the companion.\footnote{\texttt{registry-analysis/} in the source tree
of~\cite{deadly-race} gives the exact keyword patterns, per-year counts, and
caveats.} Several things stand out.

\emph{The attack is frequent and worsening.} Recorded incidents climb from a handful
a year before 2017 to a peak of $85$ in 2025 --- growth the companion asserts and
this coding counts.

\emph{The terminal move is real.} At least $5.1\%$ ($18/352$) of headlines record a
\emph{killed} or \emph{murdered} victim. A homicide is usually reported as a homicide,
not as a ``Bitcoin attack'', so this is a \emph{lower} bound --- yet even so it shows
the assassination outcome the Deadly-Race Lemma (Lemma~\ref{lem:drl}) predicts is
attested in practice, and that the cost $c$ is weighed against an action attackers
demonstrably take.

That censoring compounds with a second: obscurity does not advertise itself, so no
victim's exposure to an obscurity-dependent design is on record either. The
numerator is censored by how the terminal event is classified and the covariate by
what vendors do not say, which is why the pattern cannot be read off the data and
the bar has to be derived instead (Remark~\ref{rem:anosognosia}).

\emph{A separately sourced count agrees on the trend.} The registry is a
media-sourced convenience sample, so an independent base matters. Industry
threat-intelligence reporting --- assembled from law-enforcement disclosures,
court documents, interior-ministry statements and victim testimony rather than
from the registry --- records $72$ verified incidents in 2025, up $75\%$ year on
year, with confirmed losses of \$$40.9$M that the report itself calls an
undercount~\cite{certik-skynet-2026}; and $52$ incidents in the first half of
2026 against $39$ a year earlier, with losses of \$$124.1$M against
\$$10.5$M~\cite{certik-intel3d-h1-2026}. Stricter verification yields counts
below the registry's headline sweep, as expected; the direction, the European
concentration, and the category mix agree. Two features bear directly on this
paper. Reported \emph{home invasions} rose from $1$ to $20$ of $52$ incidents
year on year --- coercion is moving to where the holder lives, which is the
setting the model assumes. And their \emph{confirmed} homicide count is far below
our headline coding ($1$ in the half-year, against our $18$ over twelve years),
which is what two different evidentiary bars should produce and which reinforces
that the share coded here is a floor, not an estimate.

\emph{Short attacks dominate; the registry over-counts the dramatic.} Most wrench
attacks are \emph{hand-over-now} events --- a threat, an immediate transfer, minutes
on scene; the long, cinematic kidnapping is the exception. A media-sourced registry
inverts this, because a mutilation-kidnapping is news where a knifepoint mugging is
not, so the kidnap share coded below \emph{over}-states the long tail rather than
establishing it. That tacit knowledge-based authentication (TKBA,~\cite[\S2]{deadly-race})
is realistically \emph{less} effective against a
long-duration attack --- where hours of pen-and-paper effort under coercion
\emph{might} walk a deep stage --- is granted, but it is \emph{speculative} pending
the field study of the limitations section in~\cite{deadly-race}. What is \emph{not} speculative is that the
fielded alternatives fail; the cases below make that concrete.

\begin{center}\small
\begin{tabular}{@{}lrr@{}}
\toprule
\textbf{Signal (headline keyword coding)} & \textbf{Incidents} & \textbf{Share} \\
\midrule
Killed / murdered victim \emph{(a lower bound)} & $18$  & $5.1\%$  \\
Kidnap / hostage / ransom (long-hold)           & $139$ & $39.5\%$ \\
Armed robbery at gunpoint (short-hold)          & $101$ & $28.7\%$ \\
Home invasion / break-in / raid                 & $40$  & $11.4\%$ \\
Torture or physical violence                    & $82$  & $23.3\%$ \\
\bottomrule
\end{tabular}
\end{center}

\noindent Categories overlap (a kidnapping may also be violent or fatal) and do not
sum to $100\%$; shares are of all $352$ incidents, and the modality labels are a
headline proxy, not a measured duration. The sample is media-reported and
survivorship-biased --- deterred, unreported, and silently-fatal cases are absent ---
so every figure is directional and the fatality share in particular is a floor.
Remark~\ref{rem:anosognosia} gives the two mechanisms that censor it, and why they
leave the diagnosis unavailable rather than merely imprecise.

\medskip\noindent\textbf{The mix is not universal, and the Lemma's bite tracks it.} The
global shares above average over jurisdictions whose compositions differ by an order of magnitude.
Taking the ratio of long-hold to short-hold codings, $K{:}G$:

\begin{center}\small
\begin{tabular}{@{}lrrrr@{}}
\toprule
\textbf{Jurisdiction} & \textbf{$n$} & \textbf{Kidnap} & \textbf{Gunpoint} & \textbf{$K{:}G$} \\
\midrule
Brazil \emph{(small $n$)}  & $12$  & $75.0\%$ & $8.3\%$  & $9.0$ \\
France                     & $61$  & $57.4\%$ & $8.2\%$  & $7.0$ \\
\emph{All incidents}       & $352$ & $39.5\%$ & $28.7\%$ & $1.4$ \\
United States              & $59$  & $28.8\%$ & $40.7\%$ & $0.7$ \\
Russia \emph{(small $n$)}  & $10$  & $30.0\%$ & $50.0\%$ & $0.6$ \\
\bottomrule
\end{tabular}
\end{center}

\noindent The comparison that carries weight is France against the United States: the two largest
national subsets, of near-identical size ($61$ and $59$), with \emph{inverted} compositions --- a
tenfold difference in $K{:}G$. This is the sharpest available check on the media-inflation caveat
just stated. That caveat explains a \emph{level} --- why the registry over-counts dramatic incidents
everywhere --- but it cannot explain this \emph{ordering} without the further claim that French
reporting is ten times more kidnap-selective than American, in a market that is at least as heavily
covered. The parsimonious reading is that modality composition genuinely varies with local
conditions --- firearm availability, the presence of organized kidnap-for-ransom capability --- as
it does for kidnapping outside this domain entirely.

Two confounds remain and are not dismissed. Outside France and the United States the subsets are
too small to carry an argument, so the other rows are directional only. And the coding is applied to
English-language headlines, several of them translations, so a systematic vocabulary difference ---
French sources reaching for ``kidnapped'' where American ones reach for ``robbed at gunpoint'' ---
would produce part of this spread without any difference in the underlying events.

If the composition is real, it matters for scope rather than for validity. The Deadly-Race Lemma
(Lemma~\ref{lem:drl}) is a statement about the \emph{released} holder, and release is the trigger
(the release-trigger remark in~\cite{deadly-race}); a hand-over-now robbery at gunpoint approximates the clean
$p_A = 1$ endpoint --- the attacker takes what is reachable and leaves --- whereas a multi-hour hold
ending in a deliberate release is exactly the gray outcome. Coercion risk is therefore not merely
higher in some jurisdictions but \emph{differently shaped}, and the No-Gray-Area Criterion
(Criterion~\ref{prin:nga}) does its work where long holds dominate. That is the regime the record is
currently growing into: France alone accounts for a majority of the incidents recorded in 2026.

\medskip\noindent\textbf{The attacker often does not need the holder at all.} A fourth
signal falls outside the keyword table because it was coded by hand: in $12$ of the $352$
incidents ($3.4\%$) the person seized was \emph{not} a holder but a relative held to compel
one --- a parent, a spouse, an adult child.\footnote{Keyword screening surfaced $16$ candidates;
each was read individually and four were discarded as false positives, including one in which
``son of'' named a \emph{perpetrator} rather than a victim. The companion's
\texttt{registry-analysis/}~\cite{deadly-race} carries the screen and the full audited list. Because the coding reads a one-line headline, an incident that
does not name a relative is no evidence that none was involved: $3.4\%$ is a floor, and the coding
can only understate it.} Two features of this subset matter more than its size.

It is \emph{recent and concentrating}: eleven of the twelve fall in 2023 or later, and \emph{nine}
occur in a single jurisdiction (France) whose recorded incidents rose from about one a year through
2024 to $22$ in 2025 and $34$ by August 2026. Whatever share of that rise is genuine incidence
rather than reporting intensity, the tactic is not a historical scatter; it is being adopted.

It is \emph{not} the hostage token of the hostage-token property in~\cite{deadly-race}, and the difference matters. There the
leverage is a \emph{seizable object} the protocol requires, which devicelessness removes by leaving
no such object; here it is a \emph{person}, whom no custody design renders unseizable. What this
subset shows in the field is the no-delegate-coercion clause of~\cite{deadly-race} failing:
such an attacker needs no
cryptographic access and no cooperation a protocol could withhold --- only leverage, and a
relative supplies it. This bears directly on the delegated route~(\ref{r:remotewinner}), whose
clearance rests on the deciding party lying beyond $\Attacker$'s coercive reach. The delegate a
holder is most likely to actually appoint --- a spouse, a parent, an adult child --- is precisely
the party this subset shows being seized. The route is not thereby refuted: it is bounded. A
delegate who is reachable is not a second racer but a second victim, so \ref{r:remotewinner}
clears the bar only for delegates whose unreachability is structural rather than assumed --- and
that requirement, not the arrangement's convenience, is what the data prices.

\medskip\noindent\textbf{A case in point --- Balland.} The January 2025 abduction of
David Balland, a co-founder of the hardware-wallet maker Ledger~\cite{balland}, is
instructive precisely because it is an \emph{embarrassment to the self-custody creed}.
A professional champion of self-custody was seized from his home and mutilated --- a
finger severed to force payment --- and neither wealth nor prominence prevented it. He
was freed by an elite \emph{state} police unit, and the ransom --- paid not by Balland
but by an associate --- was clawed back largely because part of it moved in a
\emph{centralized} stablecoin whose issuer could freeze it. A self-custody pioneer
made whole by the state and by centralization: two mechanisms external to, and in
tension with, the creed. Only the issuer freeze is independently verifiable; the
remainder rests on the parties' accounts, according to which roughly $5\%$ was never recovered.
The case also indicts the delegated route (\ref{r:remotewinner}): an associate
\emph{did} pay, which is exactly the no-delegate-coercion clause of~\cite{deadly-race}
assumes away --- a real instance of the indirect coercion,
routed through the victim's body to a third party who will not refuse. For a brief
hand-over-now robbery a remote delegate is unreachable in time and so seems to help;
but merely \emph{not carrying} the device --- leaving the wallet in a home vault ---
does the same, while for the long attack the delegate is not out of reach but is the
ransom channel itself. The lesson is not that Balland erred: it is that, against a
determined long-hold attacker, the fielded options collapse to state rescue and
centralized freezing, whereas the tacit route removes the leverage --- no standing
device, no hostage token, no delegate to pay --- rather than relying on ex-post
recovery.

\medskip\noindent\textbf{A case in point --- RATIRL.}\label{par:ratirl} The registry coded above records
the obscurity route's failure directly. In one of its 2026 incidents (Sweden, 13
January)~\cite{ratirl} a League of Legends streamer was bound and tortured for hours by
armed intruders demanding cryptocurrency he denied holding; unable to credit the
denial, they persisted long past any value the robbery could rationally return,
leaving finally with only his Pok\'emon-card collection. This is
Corollary~\ref{cor:disclosure} in the field: where the community norm is ``deny that
you own Bitcoin,'' a coerced denial is \emph{pre-discredited} --- presumed to be the
denial-script --- so a Kerckhoffs-educated $\Attacker$ gets no exculpatory update, and
the rational continuation of an interrogation that opens with a denial is \emph{more}
coercion, not less. The case fixes the
\emph{sign}: obscurity does not merely fail to protect --- it can convert a robbery
into a torture --- \emph{negative} coercion resistance, not a mere absence. And it does so on \emph{every}
reading of the denial, which forks three ways --- each indicting obscurity.
\textbf{(i)}~If the denial was a defensive lie (there was a stash), secrecy still did
not avert the assault --- a textbook Kerckhoffs failure --- and his only residual
recourse is to \emph{double down} on it and wager against a repeat. \textbf{(ii)}~If it
was true (no stash), he bore obscurity's \emph{negative externality}: a genuine
non-holder attacked only because real holders are coached to look exactly like him ---
the very coaching that collapses the statistical holder/non-holder distinction.
\textbf{(iii)}~If the truth is in between --- a real but smaller-than-presumed stash,
too small to be worth the judicial exposure the attack incurred --- \emph{both}
pathologies operate at once: others' obscurity redirected the attack onto him, and his
own failed to stop it. That all three are simultaneously plausible is itself the
verdict: in the aggregate, across a growing population of holders and lookalikes, each
is already occurring --- and the obscurity culture, exactly as Kerckhoffs predicts,
prevents none of them.
